% Options for packages loaded elsewhere
\PassOptionsToPackage{unicode}{hyperref}
\PassOptionsToPackage{hyphens}{url}
\documentclass[
  ignorenonframetext,
]{beamer}
\newif\ifbibliography
\usepackage{pgfpages}
\setbeamertemplate{caption}[numbered]
\setbeamertemplate{caption label separator}{: }
\setbeamercolor{caption name}{fg=normal text.fg}
\beamertemplatenavigationsymbolsempty
% remove section numbering
\setbeamertemplate{part page}{
  \centering
  \begin{beamercolorbox}[sep=16pt,center]{part title}
    \usebeamerfont{part title}\insertpart\par
  \end{beamercolorbox}
}
\setbeamertemplate{section page}{
  \centering
  \begin{beamercolorbox}[sep=12pt,center]{section title}
    \usebeamerfont{section title}\insertsection\par
  \end{beamercolorbox}
}
\setbeamertemplate{subsection page}{
  \centering
  \begin{beamercolorbox}[sep=8pt,center]{subsection title}
    \usebeamerfont{subsection title}\insertsubsection\par
  \end{beamercolorbox}
}
% Prevent slide breaks in the middle of a paragraph
\widowpenalties 1 10000
\raggedbottom
\AtBeginPart{
  \frame{\partpage}
}
\AtBeginSection{
  \ifbibliography
  \else
    \frame{\sectionpage}
  \fi
}
\AtBeginSubsection{
  \frame{\subsectionpage}
}
\usepackage{iftex}
\ifPDFTeX
  \usepackage[T1]{fontenc}
  \usepackage[utf8]{inputenc}
  \usepackage{textcomp} % provide euro and other symbols
\else % if luatex or xetex
  \usepackage{unicode-math} % this also loads fontspec
  \defaultfontfeatures{Scale=MatchLowercase}
  \defaultfontfeatures[\rmfamily]{Ligatures=TeX,Scale=1}
\fi
\usepackage{lmodern}
\ifPDFTeX\else
  % xetex/luatex font selection
\fi
% Use upquote if available, for straight quotes in verbatim environments
\IfFileExists{upquote.sty}{\usepackage{upquote}}{}
\IfFileExists{microtype.sty}{% use microtype if available
  \usepackage[]{microtype}
  \UseMicrotypeSet[protrusion]{basicmath} % disable protrusion for tt fonts
}{}
\makeatletter
\@ifundefined{KOMAClassName}{% if non-KOMA class
  \IfFileExists{parskip.sty}{%
    \usepackage{parskip}
  }{% else
    \setlength{\parindent}{0pt}
    \setlength{\parskip}{6pt plus 2pt minus 1pt}}
}{% if KOMA class
  \KOMAoptions{parskip=half}}
\makeatother
\usepackage{longtable,booktabs,array}
\newcounter{none} % for unnumbered tables
\usepackage{calc} % for calculating minipage widths
\usepackage{caption}
% Make caption package work with longtable
\makeatletter
\def\fnum@table{\tablename~\thetable}
\makeatother
\setlength{\emergencystretch}{3em} % prevent overfull lines
\providecommand{\tightlist}{%
  \setlength{\itemsep}{0pt}\setlength{\parskip}{0pt}}
\usepackage{bookmark}
\IfFileExists{xurl.sty}{\usepackage{xurl}}{} % add URL line breaks if available
\urlstyle{same}
\hypersetup{
  pdftitle={Quantum Active Inference and Topological Semantic Flow},
  hidelinks,
  pdfcreator={LaTeX via pandoc}}

\title{\texorpdfstring{Quantum Active Inference and Topological Semantic
Flow}{Quantum Active Inference and Topological Semantic Flow}}
\author{\texorpdfstring{}{}}
\date{}

\begin{document}
\frame{\titlepage}

\begin{frame}{Quantum Active Inference: Topological Diagrams as Semantic
Computation}
\protect\phantomsection\label{sec:quantum-active-inference}
The preceding sections have established that categorical string
diagrams---from DisCoCat's pregroup derivations
(\autoref{sec:categorical-semantics}) through enriched hom-values
(\autoref{sec:enriched-categories}) to topos-theoretic transfer
(\autoref{sec:topos-theory})---provide a unified diagrammatic language
for case-theoretic reasoning. This section extends the framework into
its natural quantum generalization: topological quantum neural networks
(TQNNs), the ZX-calculus, and sheaf-theoretic quantum semantic
communication. The central observation is that the same
monoidal-categorical architecture that underlies DisCoCat string
diagrams also underlies quantum circuits, TQFT cobordisms, and quantum
information flow---making active inference on case-marked relational
structure a quantum topological computation.

\begin{block}{Topological Quantum Neural Networks and Spin-Network
Representations}
\protect\phantomsection\label{topological-quantum-neural-networks-and-spin-network-representations}
\begin{block}{From Quantum Neural Networks to TQFT}
\protect\phantomsection\label{from-quantum-neural-networks-to-tqft}
Marcianò, Fields, and Glazebrook demonstrate that quantum neural
networks (QNNs) admit a topological reformulation via spin-networks
{[}@fields2022tqnn{]}. In their analysis, a QNN layer or connectivity
pattern is represented as a graph whose edges carry representation
labels (spins) and whose vertices carry intertwiners---precisely the
graphical data of a spin-network in a 3-dimensional topological quantum
field theory:

\begin{quote}
``Quantum Neural Networks (QNNs) can be mapped onto spin-networks, with
the consequence that the level of analysis of their operation can be
carried out on the side of Topological Quantum Field Theory (TQFT).''
{[}@fields2022tqnn{]}
\end{quote}

This reformulation has three important structural consequences. First,
the neural architecture becomes a topological diagram (spin-network /
ribbon graph) whose evaluation by a TQFT functor gives a quantum
process; the edges and nodes of the graph encode how quantum information
flows and transforms. Second, the TQFT evaluation functorially assigns
to boundary Hilbert spaces (inputs/outputs) linear maps obtained from
topological invariants (Reshetikhin--Turaev, Turaev--Viro), playing the
role of the network's forward pass: amplitudes ``flow'' through the
topological diagram. Third, the topological control structure encodes
information flow through the network's wiring topology rather than
through any fixed geometric embedding {[}@fields2023tensor{]}.
\end{block}

\begin{block}{Universal Quantum Computation via TQNNs}
\protect\phantomsection\label{universal-quantum-computation-via-tqnns}
Fields and collaborators further demonstrate that TQNNs are universal
quantum computers by identifying the Reshetikhin--Turaev invariant of a
TQNN with a Turaev--Viro quantum error-correcting code:

\begin{quote}
``TQNNs enable universal quantum computation, using the
Reshetikhin-Turaev and Turaev-Viro models to show how TQNNs implement
quantum error-correcting codes.'' {[}@fields2025amplituhedra{]}
\end{quote}

The universality result is established via the concept of an
\emph{execution trace} for a quantum computation, leading to the
representation of TQNNs in terms of the positive geometries provided by
amplituhedra---a deep connection between quantum computation, scattering
amplitudes, and topological combinatorics.
\end{block}

\begin{block}{Quantum Reference Frames and Holographic Screens}
\protect\phantomsection\label{quantum-reference-frames-and-holographic-screens}
Fields and Glazebrook's work on quantum reference frames (QRFs) and
holographic screens provides additional algebraic structure
{[}@fields2021qrf{]}. A holographic screen---the information boundary
between two interacting quantum systems---carries a qubit array encoding
their interaction. The key insight is that QRFs deployed to identify
systems and select pointer states induce decoherence, breaking the
symmetry of the holographic encoding in an observer-relative way. This
symmetry-breaking is precisely the mechanism by which a TQNN
``observes'' its input: the choice of QRF determines the basis in which
the spin-network is evaluated.

For case-theoretic reasoning, this connects to the grammatical observer
problem: a parser or comprehender selecting a case-assignment frame for
a sentence is analogous to deploying a QRF that fixes the pointer basis
for a quantum measurement on a holographic screen.
\end{block}
\end{block}

\begin{block}{The ZX-Calculus: Categorical Quantum Circuit Diagrams}
\protect\phantomsection\label{the-zx-calculus-categorical-quantum-circuit-diagrams}
\begin{block}{String Diagrams for Quantum Processes}
\protect\phantomsection\label{string-diagrams-for-quantum-processes}
The ZX-calculus provides a categorical, diagrammatic language in which
quantum circuits are drawn as string diagrams in a symmetric monoidal
category of finite-dimensional Hilbert spaces and linear maps
{[}@kissinger2020zx{]}:

\begin{quote}
``The ZX-calculus is a graphical language for reasoning about quantum
computations and circuits\ldots{} it can represent any linear map, and
can be considered a diagrammatically complete generalization of the
usual circuit representation.'' {[}@coeckekissinger2017{]}
\end{quote}

Several structural features are critical for the connection to
case-theoretic diagrams:

\begin{itemize}
\tightlist
\item
  \textbf{Diagrams as morphisms}: ZX-diagrams are string diagrams in a
  \(\dagger\)-compact closed category. Wires represent objects (qubits),
  spiders and boxes represent morphisms, and composition/tensor product
  correspond to vertical/horizontal gluing. No metric geometry
  enters---only the topology of connections. ZX is therefore already a
  \emph{topological} representation of quantum processes.
\item
  \textbf{Circuit extraction via generalized flow}: Kissinger and van de
  Wetering demonstrate that quantum circuits can be mapped into
  ZX-diagrams, subjected to purely graph-theoretic (topological)
  transformations, and then extracted as optimized circuits. They prove
  that ``the underlying graph of our simplified ZX-diagram always has a
  graph-theoretic property called generalised flow, which in turn yields
  a deterministic circuit extraction procedure'' {[}@kissinger2020zx{]}.
\item
  \textbf{Category-theoretic semantics}: The semantics of a ZX-diagram
  are determined entirely by how components are wired
  together---precisely the compositional principle underlying DisCoCat
  and the case categories of \autoref{sec:case-systems}.
\end{itemize}
\end{block}

\begin{block}{From DisCoCat to ZX: A Shared Categorical Architecture}
\protect\phantomsection\label{from-discocat-to-zx-a-shared-categorical-architecture}
The structural parallel between pregroup grammar diagrams and
ZX-diagrams is not accidental. Both are instances of the same
mathematical object: morphisms in a compact closed monoidal category
with functorial semantics into Hilbert spaces. In DisCoCat, the functor
assigns to each grammatical type a vector space and to each derivation a
linear map computing sentence meaning {[}@coecke2010mathematical{]}. In
ZX, the standard semantics functor assigns to each spider/Hadamard
configuration a linear map in \textbf{FHilb}. The shared categorical
architecture means:

\begin{enumerate}
\tightlist
\item
  \textbf{Pregroup contractions (cups) and ZX spiders} are both
  instances of the same algebraic operation: evaluation morphisms in a
  compact closed category.
\item
  \textbf{DisCoCat normal forms} and \textbf{ZX simplifications} are
  both applications of the same rewriting theory: equational reasoning
  modulo the axioms of a compact closed category.
\item
  \textbf{The snake equation} (Cap \(\circ\) Cup = identity) that
  grounds all pregroup type reductions
  (\autoref{sec:categorical-semantics}) is a special case of the spider
  fusion rule in ZX.
\end{enumerate}

This means that case-theoretic DisCoCat derivations can, in principle,
be compiled into ZX circuits and executed on quantum hardware---a
connection already exploited by the lambeq quantum NLP pipeline
{[}@lorenz2023lambeq{]}.
\end{block}
\end{block}

\begin{block}{Aligning TQNNs with Categorical Distributional Semantics}
\protect\phantomsection\label{aligning-tqnns-with-categorical-distributional-semantics}
\begin{block}{Common Language: Ribbon and Tensor Diagrams}
\protect\phantomsection\label{common-language-ribbon-and-tensor-diagrams}
Both ZX-diagrams and TQNNs are topological string diagrams evaluated by
monoidal functors into the category of Hilbert spaces and linear maps.
The alignment becomes explicit when stated precisely:

\begin{itemize}
\tightlist
\item
  \textbf{For TQNNs}: A 3-dimensional TQFT functor from a cobordism or
  skein category to \textbf{Hilb} assigns to each spin-network/ribbon
  graph a linear map representing the TQNN computation. The underlying
  topological skein theory treats network layers as \emph{ribbon graphs}
  whose evaluation via Reshetikhin--Turaev and Turaev--Viro invariants
  gives quantum processes implementing computation and error correction
  {[}@fields2022tqnn; @fields2025amplituhedra{]}.
\item
  \textbf{For ZX}: The standard semantics functor from the free
  \(\dagger\)-compact category generated by Z/X spiders, Hadamards, etc.
  into \textbf{FHilb} assigns each ZX-diagram a linear map
  {[}@kissinger2020zx{]}.
\item
  \textbf{For DisCoCat}: The meaning functor from the pregroup grammar
  category to \textbf{FdVect} assigns to each grammatical derivation a
  multilinear map computing compositional meaning
  {[}@coecke2010mathematical{]}.
\end{itemize}

Up to choice of labels and normalization, all three are graphical
calculi for monoidal categories whose morphisms are quantum (or
quantum-like) processes. A layer of a topological quantum flow network
can be modeled as a ZX-diagram fragment whose input and output boundary
wires are the ``feature spaces'' (Hilbert spaces) at successive
processing steps, while internal spiders encode unitary/non-unitary
channels that realize synaptic transformations.
\end{block}

\begin{block}{Neural Flow as Generalized Flow}
\protect\phantomsection\label{neural-flow-as-generalized-flow}
The \emph{generalized flow} condition used to guarantee deterministic
circuit extraction from ZX-diagrams is a graph-theoretic constraint that
ensures a well-defined causal ordering of operations
{[}@kissinger2020zx{]}. This mirrors the requirement in TQNNs that the
diagram encode a consistent \emph{execution trace} of a quantum
computation. Fields and colleagues make this connection explicit in
their TQNN--amplituhedron correspondence:

\begin{quote}
``\ldots this formal correspondence is stated by Theorem 2 whose proof
draws upon the concept of execution trace for a quantum
computation\ldots{} and thus leads to representing a TQNN in terms of
the positive geometries as provided by amplituhedra.''
{[}@fields2025amplituhedra{]}
\end{quote}

A topological quantum flow neural network can therefore be regarded as a
ZX-style circuit where the graphical calculus is enriched to a 3D TQFT
skein theory, but the \emph{abstract type} of object---a topological
string diagram with functorial semantics---remains the same.
\end{block}
\end{block}

\begin{block}{Distributional Semantics on Topological Quantum Flow}
\protect\phantomsection\label{distributional-semantics-on-topological-quantum-flow}
\begin{block}{Meaning Spaces as Hilbert Spaces}
\protect\phantomsection\label{meaning-spaces-as-hilbert-spaces}
To connect TQNNs and ZX circuits to \emph{distributional semantics}, one
reinterprets the amplitudes and correlations in these diagrams as
semantic quantities. Recent work on quantum semantic communication
provides the theoretical bridge, modeling ``meaning spaces'' as Hilbert
spaces at nodes of a graph with channels as completely positive
trace-preserving (CPTP) maps along edges {[}@khatri2025quantum{]}:

\begin{quote}
``Multi-agent semantic networks are modeled as quantum sheaves, where
agents' meaning spaces are Hilbert spaces connected by quantum
channels.'' {[}@khatri2025quantum{]}
\end{quote}

A \emph{quantum semantic sheaf} over a communication graph \(G=(V,E)\)
is defined as a triple \((H,F,\rho)\) where each vertex \(v\) carries a
finite-dimensional Hilbert space \(H_v\), edges carry CPTP maps \(F_e\),
and each vertex has a local density operator \(\rho_v\) encoding its
``current semantic state.'' This is precisely a distributional-semantics
picture: meanings are modeled as vectors/density operators in
high-dimensional spaces, with co-occurrence and context encoded in how
these spaces are functorially related across the network.
\end{block}

\begin{block}{Grafting Distributional Semantics onto the TQNN/ZX
Architecture}
\protect\phantomsection\label{grafting-distributional-semantics-onto-the-tqnnzx-architecture}
When this sheaf-theoretic semantics is grafted onto the TQNN/ZX
architecture, three structural correspondences emerge:

\begin{enumerate}
\item
  \textbf{Edges as semantic feature channels}: In a TQNN or ZX-diagram,
  each wire carries not merely an abstract qubit but a semantic Hilbert
  space \(H_v\) associated with a context, concept, or agent. Amplitudes
  or density matrices on that wire encode a distribution over semantic
  features---exactly as word vectors encode distributional meaning in
  classical compositional distributional semantics.
\item
  \textbf{Nodes as compositional operations}: Spiders/gates in ZX or
  intertwiners in TQNNs become \emph{semantic composition maps}: they
  take distributed meanings on input wires and produce new distributed
  meanings on output wires, analogous to how DisCoCat composes word
  meanings into phrase/sentence meanings via multilinear maps.
\item
  \textbf{Topological wiring as contextual structure}: The topology of
  the diagram---the way wires and nodes are connected---encodes which
  semantic spaces interact and in what causal/structural pattern. This
  is the semantic analogue of syntactic structure in distributional
  semantics, realized as a topological quantum circuit.
\end{enumerate}

In this reading, a topological quantum flow neural network becomes a
\emph{distributional semantic machine}: a functor that sends a
topological diagram (graph of contexts and interactions) to a family of
Hilbert spaces and maps where vectors/densities represent distributed
meanings and their probabilistic transformations.
\end{block}
\end{block}

\begin{block}{Semantic Probabilistic Information Transfer as Functorial
Flow}
\protect\phantomsection\label{semantic-probabilistic-information-transfer-as-functorial-flow}
\begin{block}{Sheaf Cohomology and Semantic Alignment}
\protect\phantomsection\label{sheaf-cohomology-and-semantic-alignment}
The sheaf-based framework proves that semantic alignment between agents
is governed by cohomology classes of the quantum semantic sheaf;
contextuality and entanglement act as resources that remove obstructions
to alignment {[}@khatri2025quantum{]}:

\begin{quote}
``We derive semantic channel capacity when sender and receiver share
prior entanglement, proving it strictly exceeds classical capacity. The
quantum advantage grows as channel noise increases---precisely when
semantic communication most benefits over bit-level transmission.''
{[}@khatri2025quantum{]}
\end{quote}

Two further results from this framework are particularly relevant:

\begin{itemize}
\tightlist
\item
  \textbf{Contextuality as a semantic resource}: ``Quantum contextuality
  reduces cohomological obstructions to semantic alignment. Contextual
  correlations act as `pre-shared semantic resolution,' establishing
  contextuality as a resource for semantic communication''
  {[}@khatri2025quantum{]}.
\item
  \textbf{Discord as integrated semantic information}: ``Quantum discord
  equals integrated semantic information, linking quantum correlations
  to irreducible semantic content and connecting our framework to
  integrated information theory'' {[}@khatri2025quantum{]}.
\end{itemize}

These results establish that the topology of the semantic sheaf (and its
cohomology) constrains how probabilistic semantic information can be
transferred; quantum features (entanglement, contextuality, discord)
change these constraints in well-defined ways.
\end{block}

\begin{block}{Topological Circuit as Semantic Sheaf Skeleton}
\protect\phantomsection\label{topological-circuit-as-semantic-sheaf-skeleton}
A TQNN/ZX circuit implementing a quantum communication or computation
protocol is itself a diagram over which one can define a sheaf of
semantic spaces and channels. The underlying graph of the TQNN/ZX
diagram serves as the base graph \(G=(V,E)\) of the semantic sheaf:
vertices inherit meaning spaces \(H_v\), edges inherit CPTP maps
\(F_e\), and the TQFT/ZX functor gives the global linear map
representing the protocol. Distributional semantics as diagram
evaluation then becomes literal: passing an initial semantic state
(distribution over meanings) through the TQNN/ZX diagram yields an
output state whose components encode the posterior semantic
distributions at boundary wires.

ZX rewrite rules, which change the internal topology of the diagram
while preserving its overall semantics as a linear map, correspond to
alternative factorizations of the same semantic
transformation---different internal ``flow architectures'' for realizing
the same semantic map.
\end{block}
\end{block}

\begin{block}{Implications for Active Inference and Case Assignment}
\protect\phantomsection\label{implications-for-active-inference-and-case-assignment}
\begin{block}{Case Assignment as Quantum Measurement}
\protect\phantomsection\label{case-assignment-as-quantum-measurement}
The active inference model of case reasoning
(\autoref{sec:cognitive-integration}) acquires a new dimension in this
quantum topological setting. Case assignment---the cognitive process of
determining \emph{who does what to whom}---can be modeled as a quantum
measurement process on a holographic screen, with the following
correspondences:

{\def\LTcaptype{none} % do not increment counter
\begin{longtable}[]{@{}
  >{\raggedright\arraybackslash}p{(\linewidth - 2\tabcolsep) * \real{0.5000}}
  >{\raggedright\arraybackslash}p{(\linewidth - 2\tabcolsep) * \real{0.5000}}@{}}
\toprule\noalign{}
\begin{minipage}[b]{\linewidth}\raggedright
Classical Case Assignment
\end{minipage} & \begin{minipage}[b]{\linewidth}\raggedright
Quantum Topological Model
\end{minipage} \\
\midrule\noalign{}
\endhead
Case role (NOM, ACC, \ldots) & Pointer state selected by QRF \\
Case frame (alignment system) & Quantum reference frame \\
Relational structure of event & Spin-network topology \\
Free-energy minimization & TQFT evaluation of diagram \\
Prediction-error (P600/N400) & Symmetry-breaking on holographic
screen \\
\bottomrule\noalign{}
\end{longtable}
}
\end{block}

\begin{block}{From Predictive Processing to Topological Flow}
\protect\phantomsection\label{from-predictive-processing-to-topological-flow}
In the predictive processing account, a cognitive agent maintains a
generative model that predicts the relational structure of incoming
linguistic material. When this model is realized as a TQNN, prediction
becomes evaluation of the topological diagram; prediction error becomes
the discrepancy between the predicted TQFT evaluation and the observed
data; and belief updating becomes modification of the spin-network's
edge labels (representation labels) and vertex intertwiners.

The topological character of this computation confers significant
advantages for active inference on case structure: topological
invariants are robust to continuous deformation, so the generative
model's predictions are stable under small perturbations of the
input---a desirable property for language understanding in noisy
environments.
\end{block}

\begin{block}{Quantum Advantage in Semantic Communication}
\protect\phantomsection\label{quantum-advantage-in-semantic-communication}
The sheaf-theoretic results of Khatri et al. {[}-@khatri2025quantum{]}
suggest that quantum features provide genuine advantages for semantic
communication---not merely computational speedup, but qualitative
enhancements in semantic alignment. If case-marked relational structure
is communicated between agents via quantum channels, entanglement
provides additional semantic capacity, contextuality removes alignment
obstructions, and discord captures irreducible semantic content. These
are not abstract possibilities but operational consequences of the
mathematical framework developed across this review.
\end{block}
\end{block}
\end{frame}

\end{document}
